% GENERATED FILE. Edit canonical lessons, then run npm run generate:books.
% canonical-source-hash: fnv1a64:6ef8a9186036b465
% canonical-lessons: PA-C263-mor, PA-C263-murna, PA-C263-pasa, PA-C263-sajja, PA-C263-khabba

\chapter{Right, Left and the Bend}
\label{ch:pa-a2-right-left-and-the-bend}
\hlchaptermodality{263}

\begin{chapteropening}
\textbf{By the end of this chapter:} \emph{I can say which way to turn.}

Complete the last lesson of chapter 263: I can say which way to turn.
\end{chapteropening}

\section[\texorpdfstring{\pa{ਮੋੜ} (moṛ)}{ਮੋੜ (moṛ)}]{\pa{ਮੋੜ} (moṛ) — a bend in the road; a turning}

\label{lesson:PA-C263-mor}



\begin{quote}
Before the new one: say the Punjabi for a receipt, then the Punjabi for to pay off.
\end{quote}

\subsection*{You'll want to know: \pa{ਮੋੜ}}
\textbf{\pa{ਮੋੜ}} — \emph{moṛ} — \textquotedblleft{}a bend in the road; a turning\textquotedblright{}.

\begin{quote}
\pa{ਅਗਲੇ} \pa{ਮੋੜ} \pa{ਉੱਤੇ} \pa{ਇੱਕ} \pa{ਦੁਕਾਨ} \pa{ਹੈ।}
\end{quote}

\begin{quote}
\emph{agle moṛ utte ikk dukān hai.}
\end{quote}

\begin{quote}
There is a shop at the next bend.
\end{quote}

A masculine noun.

\begin{cousinweb}
\textbf{\pa{ਮੋੜ}} (\emph{moṛ}) is inherited from Sanskrit \emph{muṭati}, \textquotedblleft{}twists\textquotedblright{}; CDIAL lists Punjabi \emph{moṛ}, \textquotedblleft{}turn in a road\textquotedblright{}, under it (CDIAL 10186).
\end{cousinweb}

\begin{grammarlens}[title={What you've built}]
One more word for this chapter.
\end{grammarlens}

\subsection*{Your turn}
\begin{itemize}
\raggedright
  \item \emph{Say it:} \emph{moṛ}
  \item \emph{Say it:} \emph{moṛ}, once more
  \item \emph{Say it:} \emph{agle moṛ utte ikk dukān hai.} — There is a shop at the next bend.
  \item \emph{Recall:} say the Punjabi for a receipt, then the Punjabi for to pay off, then say \emph{moṛ} again
\end{itemize}

\begin{tcolorbox}[breakable,colback=teal!4,colframe=teal!35!black,title={Before you move on}]
What does \pa{ਮੋੜ} mean? (\textquotedblleft{}A bend in the road; a turning\textquotedblright{}.) Say it once more.
\end{tcolorbox}
\section[\texorpdfstring{\pa{ਮੁੜਨਾ} (muṛnā)}{ਮੁੜਨਾ (muṛnā)}]{\pa{ਮੁੜਨਾ} (muṛnā) — to turn; to turn off}

\label{lesson:PA-C263-murna}



\begin{quote}
Before the new one: say the Punjabi for to pay off, then the Punjabi for a bend in the road.
\end{quote}

\subsection*{You'll want to know: \pa{ਮੁੜਨਾ}}
\textbf{\pa{ਮੁੜਨਾ}} — \emph{muṛnā} — \textquotedblleft{}to turn; to turn off\textquotedblright{}.

\begin{quote}
\pa{ਬੱਸ} \pa{ਚੌਂਕ} \pa{ਤੋਂ} \pa{ਮੁੜਦੀ} \pa{ਹੈ।}
\end{quote}

\begin{quote}
\emph{bass chauṅk ton muṛdī hai.}
\end{quote}

\begin{quote}
The bus turns at the square.
\end{quote}

Past (I …): \pa{ਮੈਂ} \pa{ਮੋੜ} \pa{ਤੋਂ} \pa{ਮੁੜਿਆ}, \emph{maiṁ moṛ ton muṛiā}.

\begin{cousinweb}
\textbf{\pa{ਮੁੜਨਾ}} (\emph{muṛnā}) is inherited from Sanskrit \emph{muṭati}, \textquotedblleft{}twists\textquotedblright{}; CDIAL lists Punjabi \emph{muṛṇā}, \textquotedblleft{}to be twisted\textquotedblright{}, under it (CDIAL 10186).
\end{cousinweb}

\begin{grammarlens}[title={What you've built}]
One more word for this chapter.
\end{grammarlens}

\subsection*{Your turn}
\begin{itemize}
\raggedright
  \item \emph{Say it:} \emph{muṛnā}
  \item \emph{Say it:} \emph{muṛnā}, once more
  \item \emph{Say it:} \emph{bass chauṅk ton muṛdī hai.} — The bus turns at the square.
  \item \emph{Recall:} say the Punjabi for to pay off, then the Punjabi for a bend in the road, then say \emph{muṛnā} again
  \item \emph{Write it:} copy \textbf{\pa{ਮੁੜਨਾ}} once, checking the whole word
\end{itemize}

\begin{tcolorbox}[breakable,colback=teal!4,colframe=teal!35!black,title={Before you move on}]
What does \pa{ਮੁੜਨਾ} mean? (\textquotedblleft{}To turn; to turn off\textquotedblright{}.) Say it once more.
\end{tcolorbox}
\section[\texorpdfstring{\pa{ਪਾਸਾ} (pāsā)}{ਪਾਸਾ (pāsā)}]{\pa{ਪਾਸਾ} (pāsā) — a side}

\label{lesson:PA-C263-pasa}



\begin{quote}
Before the new one: say the Punjabi for a bend in the road, then the Punjabi for to turn.
\end{quote}

\subsection*{You'll want to know: \pa{ਪਾਸਾ}}
\textbf{\pa{ਪਾਸਾ}} — \emph{pāsā} — \textquotedblleft{}a side\textquotedblright{}.

\begin{quote}
\pa{ਸਕੂਲ} \pa{ਸੜਕ} \pa{ਦੇ} \pa{ਦੂਜੇ} \pa{ਪਾਸੇ} \pa{ਹੈ।}
\end{quote}

\begin{quote}
\emph{sakūl saṛak de dūje pāse hai.}
\end{quote}

\begin{quote}
The school is on the other side of the road.
\end{quote}

A masculine noun.

\begin{cousinweb}
\textbf{\pa{ਪਾਸਾ}} (\emph{pāsā}) is inherited from Sanskrit \emph{pārśva}, \textquotedblleft{}region of the ribs, side\textquotedblright{}; CDIAL lists Punjabi \emph{pāssā}, \textquotedblleft{}side, direction\textquotedblright{}, under it (CDIAL 8118).
\end{cousinweb}

\begin{grammarlens}[title={What you've built}]
One more word for this chapter.
\end{grammarlens}

\subsection*{Your turn}
\begin{itemize}
\raggedright
  \item \emph{Say it:} \emph{pāsā}
  \item \emph{Say it:} \emph{pāsā}, once more
  \item \emph{Say it:} \emph{sakūl saṛak de dūje pāse hai.} — The school is on the other side of the road.
  \item \emph{Recall:} say the Punjabi for a bend in the road, then the Punjabi for to turn, then say \emph{pāsā} again
  \item \emph{Write it:} copy \textbf{\pa{ਪਾਸਾ}} once, checking the whole word
\end{itemize}

\begin{tcolorbox}[breakable,colback=teal!4,colframe=teal!35!black,title={Before you move on}]
What does \pa{ਪਾਸਾ} mean? (\textquotedblleft{}A side\textquotedblright{}.) Say it once more.
\end{tcolorbox}
\section[\texorpdfstring{\pa{ਸੱਜਾ} (sajjā)}{ਸੱਜਾ (sajjā)}]{\pa{ਸੱਜਾ} (sajjā) — right-hand; on the right}

\label{lesson:PA-C263-sajja}



\begin{quote}
Before the new one: say the Punjabi for to turn, then the Punjabi for a side.
\end{quote}

\subsection*{You'll want to know: \pa{ਸੱਜਾ}}
\textbf{\pa{ਸੱਜਾ}} — \emph{sajjā} — \textquotedblleft{}right-hand; on the right\textquotedblright{}.

\begin{quote}
\pa{ਬੈਂਕ} \pa{ਸੜਕ} \pa{ਦੇ} \pa{ਸੱਜੇ} \pa{ਪਾਸੇ} \pa{ਹੈ।}
\end{quote}

\begin{quote}
\emph{baiṅk saṛak de sajje pāse hai.}
\end{quote}

\begin{quote}
The bank is on the right-hand side of the road.
\end{quote}

\begin{cousinweb}
\textbf{\pa{ਸੱਜਾ}} (\emph{sajjā}) is inherited from Sanskrit \emph{sajya}, \textquotedblleft{}strung (of a bow)\textquotedblright{}; CDIAL lists Punjabi \emph{sajjā}, \textquotedblleft{}righthand\textquotedblright{}, under it (CDIAL 13095).
\end{cousinweb}

\begin{grammarlens}[title={What you've built}]
One more word for this chapter.
\end{grammarlens}

\subsection*{Your turn}
\begin{itemize}
\raggedright
  \item \emph{Say it:} \emph{sajjā}
  \item \emph{Say it:} \emph{sajjā}, once more
  \item \emph{Say it:} \emph{baiṅk saṛak de sajje pāse hai.} — The bank is on the right-hand side of the road.
  \item \emph{Recall:} say the Punjabi for to turn, then the Punjabi for a side, then say \emph{sajjā} again
  \item \emph{Write it:} copy \textbf{\pa{ਸੱਜਾ}} once, checking the whole word
\end{itemize}

\begin{tcolorbox}[breakable,colback=teal!4,colframe=teal!35!black,title={Before you move on}]
What does \pa{ਸੱਜਾ} mean? (\textquotedblleft{}Right-hand; on the right\textquotedblright{}.) Say it once more.
\end{tcolorbox}
\section[\texorpdfstring{\pa{ਖੱਬਾ} (khabbā)}{ਖੱਬਾ (khabbā)}]{\pa{ਖੱਬਾ} (khabbā) — left-hand; on the left}

\label{lesson:PA-C263-khabba}



\begin{quote}
Before the new one: say the Punjabi for a side, then the Punjabi for right-hand.
\end{quote}

\subsection*{You'll want to know: \pa{ਖੱਬਾ}}
\textbf{\pa{ਖੱਬਾ}} — \emph{khabbā} — \textquotedblleft{}left-hand; on the left\textquotedblright{}.

\begin{quote}
\pa{ਮੋੜ} \pa{ਤੋਂ} \pa{ਖੱਬੇ} \pa{ਪਾਸੇ} \pa{ਮੁੜੋ।}
\end{quote}

\begin{quote}
\emph{moṛ ton khabbe pāse muṛo.}
\end{quote}

\begin{quote}
At the bend, turn to the left.
\end{quote}

\begin{cousinweb}
\textbf{\pa{ਖੱਬਾ}} (\emph{khabbā}) is perhaps inherited from Sanskrit \emph{kharva}, \textquotedblleft{}mutilated, imperfect\textquotedblright{}; CDIAL lists Punjabi \emph{khabbā}, \textquotedblleft{}mutilated, imperfect\textquotedblright{}, and a Sindhi form meaning \textquotedblleft{}left-handed\textquotedblright{} under it (CDIAL 3832). How the sense \textquotedblleft{}left\textquotedblright{} arose is not settled.
\end{cousinweb}

\begin{grammarlens}[title={What you've built}]
One more word for this chapter.
\end{grammarlens}

\subsection*{Your turn}
\begin{itemize}
\raggedright
  \item \emph{Say it:} \emph{khabbā}
  \item \emph{Say it:} \emph{khabbā}, once more
  \item \emph{Say it:} \emph{moṛ ton khabbe pāse muṛo.} — At the bend, turn to the left.
  \item \emph{Recall:} say the Punjabi for a side, then the Punjabi for right-hand, then say \emph{khabbā} again
  \item \emph{Write it:} copy \textbf{\pa{ਖੱਬਾ}} once, checking the whole word
\end{itemize}

\begin{tcolorbox}[breakable,colback=teal!4,colframe=teal!35!black,title={Before you move on}]
What does \pa{ਖੱਬਾ} mean? (\textquotedblleft{}Left-hand; on the left\textquotedblright{}.) Say it once more.
\end{tcolorbox}
